%======================================================================
\subsection{BW-SIM-04 --- Built, 2 Nov: the CalculiX deck writes fast on multi-solid models}\label{sec:sim-deck-built}
%======================================================================

The turbojet session spent most of a 773\,s run writing the \file{.inp} for a three-solid rotor (146k nodes, second
order, an M8 thread). The profile says where: 99.97\,\% of the deck write was \code{FemMesh.getNodesBySolid}, an OCC
point-in-solid test (\code{BRepExtrema\_DistShapeShape}, solid treatment) for every node and every solid. The deck is now
written from the mesh itself and is the same deck.

\paragraph{What you can do now.}
\begin{itemize}[nosep]
  \item Nothing changes in the panels. A multi-solid study (a material per solid) writes its input in minutes
        instead of tens of minutes; spin and centrifugal loads on a solid use the same shortcut.
  \item If the shortcut cannot decide (hex or wedge elements, a reference that is not a solid, overlapping solids, a mesh
        without surface elements) the old node search runs, unchanged.
\end{itemize}

\paragraph{How it works.}
The volume mesh keeps one surface triangle on every face of the CAD shape, including the faces two solids share.
Cutting the tetrahedra at those triangles leaves one connected piece per solid (\code{meshtools.\_MeshPieces}, scipy
connected components). Seven element centroids per piece decide which solid it is. Only the mid-edge nodes that can lie
off a solid (both end nodes on a surface triangle) are measured, and only against their own solid, as the node search
did; elements with such a node go to the existing leftover handling (nearest solid, \code{LAST\_UNREFERENCED}), so the
warning in the run pane is unchanged. Mid nodes the mesher placed on a curved surface are measured on a plane, cylinder,
cone, sphere or torus face first (microseconds) and on the solid only when that fails.

\paragraph{Decided while building.}
\begin{itemize}[nosep]
  \item \textbf{Volume ids do not survive.} A FemMesh exposes no shape ids, and the study rebuilds the mesh (orphans,
        Tet10 upgrade, folded-element repair) before the deck is written, so reading element$\to$solid from the mesher was
        not possible; the pieces come from the mesh topology.
  \item \textbf{Node sets use it too.} \code{get\_femnodes\_by\_refshape} returns the nodes of a solid's pieces without
        its off-solid mid nodes, in sorted order like the C++ set.
  \item \textbf{Freeform curved mid nodes.} Up to 2000 per solid, and any far from the chord midpoint, are all
        measured; beyond that a spread sample of 60 is measured and, if one is off, all. This is the one place where the
        result rests on the mesher having put the node on the surface; the test with a moved node shows a displaced node
        is still found.
  \item \textbf{Switch.} \code{meshtools.FAST\_SOLID\_SETS = False} restores the old search (the test uses it).
  \item \textbf{No fork pool.} Measuring nodes in forked workers deadlocks in OCC after \code{fork}; the remaining time is
        the serial OCC distance of the straight mid nodes that are not on an analytic face.
\end{itemize}

\paragraph{Timings} (repo build, FreeCADCmd; the machine was shared, so read the ratios).
\begin{center}
\begin{tabular}{lrrr}
\hline
Model (3 solids, Tet10) & old & new & deck \\
\hline
Netgen, 65k nodes, M8-like thread & 480\,s & 56\,s & identical \\
Gmsh, 59k nodes & 120\,s & 45\,s & identical (1104 leftover elements) \\
Netgen, 160k nodes & hours (starved) & 164\,s & not compared \\
Gmsh, 181k nodes & --- & 89\,s & not compared \\
\hline
\end{tabular}
\end{center}
Deck equality means the \file{Mesh.inp} files equal line for line apart from the comment lines (date, file name).

\paragraph{Tests.} \file{TestBruceWorxDeckWrite}: element sets, node sets, a pushed-off mid node (leftover elements) and
the whole \file{.inp}, old against new, on a shaft, disc and cap compound. SimCore, SimRotor, SimModel and Validation
suites pass (4, 63, 68 and 50, 32 tests); \file{TestMaterialSetsOutsideSolids} caught a corner node pushed off a face, so surface corner nodes are measured too.

\paragraph{Open.}
\begin{itemize}[nosep]
  \item The straight mid nodes still cost about 25\,ms each on a thread (one OCC solid test); a cheaper inside test would
        make the 160k-node case well under a minute.
  \item The node search itself (\code{getNodesBySolid}) is untouched; constraints on faces and edges still use it.
\end{itemize}
